\clearpage\thispagestyle{empty}
\pdfbookmark[0]{Quellen zu Heller Fleck im Mond}{moonmarch1821-sources}
\begin{center}\textbf{Quellen zu Heller Fleck im Mond}\end{center}
\noindent Diese KI-Übertragung ins heutige Deutsch folgt dem Bericht [\textit{Heller Fleck im Mond.}] in \href{https://gdz.sub.uni-goettingen.de/id/PPN236007467}{Gauß’ \textit{Werken}, Bd.~VI} (Göttingen, 1874), S.~436--437. Der Bericht zitiert aus dem Brief von Dr. Heinrich Wilhelm Olbers an Gauß vom 27. Februar 1821 einschließlich der Tagebuchaufzeichnung und verweist auf einen späteren Brief von Olbers. Der Erstdruck erschien unter Göttingen in den \textit{Göttingischen gelehrten Anzeigen}, Nr.~46, 22. März 1821, S.~\href{https://images.p.sub.uni-goettingen.de/iiif/image/gdz:PPN319721507_1821:00000451/full/1600,/0/default.jpg}{449}, \href{https://images.p.sub.uni-goettingen.de/iiif/image/gdz:PPN319721507_1821:00000452/full/1600,/0/default.jpg}{450}, \href{https://images.p.sub.uni-goettingen.de/iiif/image/gdz:PPN319721507_1821:00000453/full/1600,/0/default.jpg}{451} und \href{https://images.p.sub.uni-goettingen.de/iiif/image/gdz:PPN319721507_1821:00000454/full/1600,/0/default.jpg}{452}. Das digitalisierte Exemplar stammt aus der Niedersächsischen Staats- und Universitätsbibliothek Göttingen.
\par\smallskip\noindent Beim Vergleich mit einem Fixstern nordöstlich vom Mond nennt die Werkausgabe die dritte Größe, der Erstdruck auf S.~\href{https://images.p.sub.uni-goettingen.de/iiif/image/gdz:PPN319721507_1821:00000452/full/1600,/0/default.jpg}{450} die sechste. Die spätere hypothetische Helligkeit sechster Größe steht in beiden Fassungen. Die Werkausgabe hat \textit{Maeilius}. Auf S.~\href{https://images.p.sub.uni-goettingen.de/iiif/image/gdz:PPN319721507_1821:00000454/full/1600,/0/default.jpg}{452} des Erstdrucks steht offenbar \textit{Manilius}. Die Übertragung folgt der Werkausgabe. KI-Übertragung und KI-Satz.
\clearpage
